\section{Conclusion}
\label{sec:conclusion}

In this work, we introduced Norm-Equalized Task Arithmetic (NETA), a training-free and parameter-free closed-form method for multi-task model merging. Guided by the principle of Occam's razor, NETA addresses the challenge of task dominance in standard Task Arithmetic by analytically equalizing the Frobenius norms of task vectors at each layer before merging. This simple weight-space transformation guarantees perfect isotropic magnitude balancing across tasks at every level of representation, without introducing any new hyperparameters or requiring test-time calibration data.

Through comprehensive evaluations on CLIP ViT-B/32 across four diverse visual classification datasets, we showed that NETA outperforms vanilla Task Arithmetic and TIES-Merging across MNIST and FashionMNIST zero-shot. Crucially, our investigations exposed the \textbf{Overfitting-Optimizer Paradox} of contemporary test-time adaptation methods (such as AdaMerging). We showed that optimizing weight coefficients via joint entropy minimization on small calibration sets creates a transductive overfit state that favors easy, low-entropy tasks while severely suppressing and neglecting harder, high-entropy tasks. NETA avoids this failure mode entirely, achieving exceptional robustness and highly stable performance across different random seeds.

\subsection{Limitations and Future Work}
While NETA represents a significant step toward training-free model merging, it has certain limitations that present exciting avenues for future research. First, although we evaluated our models on standard classification task vectors, our empirical validation was limited to a moderately sized CLIP ViT-B/32 visual encoder and a 4-dataset visual suite. A primary limitation of this evaluation scope is that standard CLIP model merging benchmarks are typically conducted across an 8-dataset suite representing a wider variety of specialized domain shifts. Scaling NETA's isotropic norm-equalization to the full 8-dataset suite and to larger modern architectures---such as Large Language Models (LLMs), vision-language backbones, or generative multi-modal networks---is a crucial next step to fully establish its generalizability. Second, our analysis was restricted to a standardized subset of 1024 test images per dataset to manage computational overhead on CPU nodes under strict Slurm queue limits. While 1024 images provides a robust statistical baseline that reliably preserves performance ranks, validating NETA on full-scale, massive benchmark datasets will be valuable to confirm generalizability. Finally, while our layer-wise norm balancing achieves excellent multi-task representation fairness, it assumes that task experts should contribute equally at all stages of the network. However, early layers in deep networks typically encode general visual abstractions, whereas deeper layers encode highly specialized, task-specific features. Thus, exploring anisotropic scaling formulations---where early layers enforce strict isotropic magnitude balancing to preserve visual stream consistency while deeper layers allow selective task dominance based on task difficulty or semantic alignment---presents an extremely promising path to further optimize multi-task performance. Furthermore, extending these concepts to incorporate analytical proxies of parameter importance, such as parameter covariance or Fisher information, represents another valuable direction for future training-free model merging research.

By demonstrating that a simple, three-line analytical transformation can achieve balanced, highly competitive multi-task performance without any training, optimization parameters, or calibration data, NETA highlights the immense power of simplicity in modern deep learning. We hope this work inspires the machine learning community to seek clean, physically-grounded parameter-space solutions, rather than defaulting to complex, parameter-heavy optimization and search pipelines.
